% --------------- ABSTRACT ---------------
\chapter*{Abstract}
\addcontentsline{toc}{chapter}{Abstract}

Convolutional neural networks (CNNs) for automated visual inspection in Industry~4.0 can
base their decisions on the background instead of the workpiece. In the reviewed
literature, studies that vary the share of background use at most one saliency method,
and none compares explainability methods against a complex industrial background as that
share grows. This thesis makes that comparison on 15 cylindrical workpieces of different
materials and colours, one physical piece each, photographed in the Fischertechnik
Industry~4.0 learning factory. Crops at five object-to-image (O2I) ratios from 90\% to
10\% were classified by 75 one-vs-rest CNNs trained from scratch. Occlusion sensitivity
and Gradient-weighted Class Activation Mapping (Grad-CAM) were compared through a
per-pixel centre measure and checked against direct tests that remove the workpiece or
the background. As the ratio falls, the mean test F1 drops from 0.813 to 0.720, mostly
below 50\%, but the background does not take over the decision at any ratio: without
the workpiece the models cannot tell the materials apart, without the background they
rank the crops as well as before, and the centre measure obtained by removal stays well
above the no-preference level of 0.5. The errors are mostly confusions between
look-alike workpieces, above all aluminium, steel and grey plastic. The two methods
disagree. On average over the materials, occlusion sensitivity agrees with the direct
tests. The Grad-CAM centre
measure falls from 0.737 to 0.407, a shift that the tests do not confirm; it arises from
the averaging of gradients in Grad-CAM, and HiResCAM, which avoids it, agrees with
occlusion. A trend in a heatmap should therefore be checked by intervening on the input
before it is read as a property of the model.

\vspace{1cm}
\noindent\textbf{Keywords:} material classification, object-to-image ratio, complex background,
occlusion sensitivity, Grad-CAM, explainability, learning factory

% --------------- ZUSAMMENFASSUNG (German abstract) ---------------
\chapter*{Zusammenfassung}
\addcontentsline{toc}{chapter}{Zusammenfassung}

\begin{otherlanguage}{ngerman}
Convolutional Neural Networks (CNNs) für die automatisierte Sichtprüfung in der
Industrie~4.0 können ihre Entscheidungen auf den Hintergrund statt auf das Werkstück
stützen. In der gesichteten Literatur wird bei variiertem Hintergrundanteil höchstens
eine Saliency-Methode eingesetzt, und es liegt kein Vergleich von
Erklärbarkeitsmethoden vor einem komplexen industriellen Hintergrund bei wachsendem
Anteil vor. Diese Arbeit vergleicht zwei solche Methoden an 15 zylindrischen Werkstücken
unterschiedlicher Materialien und Farben, je ein physisches Exemplar, die in der
Fischertechnik-Industrie-4.0-Lernfabrik aufgenommen wurden. Ausschnitte mit fünf
Objekt-Bild-Verhältnissen (object-to-image, O2I) von 90\,\% bis 10\,\% wurden mit 75 von
Grund auf trainierten One-vs-Rest-CNNs klassifiziert. Okklusionssensitivität und
Gradient-weighted Class Activation Mapping (Grad-CAM) wurden über ein pixelbezogenes
Zentrumsmaß verglichen und mit direkten Tests überprüft, bei denen das Werkstück oder
der Hintergrund entfernt wird. Mit sinkendem Verhältnis fällt der mittlere
Test-F1-Wert von 0,813 auf 0,720, größtenteils unterhalb von 50\,\%, doch bei keinem
Verhältnis übernimmt der Hintergrund die Entscheidung: Ohne das Werkstück können die
Modelle die Materialien nicht mehr unterscheiden, ohne den Hintergrund ordnen sie die
Ausschnitte so gut wie zuvor, und das durch Entfernen gewonnene Zentrumsmaß bleibt
deutlich über dem neutralen Wert 0,5. Die Fehler sind überwiegend Verwechslungen
ähnlich aussehender Werkstücke, vor allem von Aluminium, Stahl und grauem Kunststoff.
Die beiden Methoden widersprechen sich. Im Mittel über die Materialien stimmt die
Okklusionssensitivität mit den direkten Tests überein. Das Grad-CAM-Zentrumsmaß fällt von 0,737 auf 0,407, eine
Verschiebung, die die Tests nicht bestätigen; sie entsteht durch die Mittelung der
Gradienten in Grad-CAM, und HiResCAM, das diese Mittelung vermeidet, stimmt mit der
Okklusion überein. Ein Trend in einer Heatmap sollte daher durch einen Eingriff in die
Eingabe geprüft werden, bevor er als Eigenschaft des Modells gelesen wird.

\vspace{1cm}
\noindent\textbf{Schlagwörter:} Materialklassifikation, Objekt-Bild-Verhältnis, komplexer
Hintergrund, Okklusionssensitivität, Grad-CAM, Erklärbarkeit, Lernfabrik
\end{otherlanguage}
